\begin{frame}{Faithful does not imply understandable}
\PerspectiveTag{\ContentIcon}{Content}\hfill \PerspectiveTag{\PresentationIcon}{Presentation}

\vspace{.4em}
\[
\text{Faithful}\not\Rightarrow\text{Understandable}
\]
A technically faithful explanation with a large number of predicates may be unusable for its audience.
\end{frame}

\begin{frame}{Plausible does not imply faithful}
\PerspectiveTag{\UserIcon}{User}\hfill \PerspectiveTag{\ContentIcon}{Content}

\vspace{.4em}
\[
\text{Coherent with human knowledge}\not\Rightarrow\text{Correct relative to model}
\]
A domain expert may find an explanation plausible even if the model relies on a spurious feature.
\end{frame}

\begin{frame}{Human-grounded tasks}
\PerspectiveTag{\UserIcon}{User}\hfill Coherence, context, controllability

\vspace{.4em}
\begin{itemize}
  \item Predict the model output; detect a model failure.
  \item Identify influential variables; answer comprehension questions.
  \item Measure decision time, confidence, and calibration.
\end{itemize}
Preference alone is weak evidence: ``Which explanation do you like?'' does not test task performance.
\end{frame}

\begin{frame}{Application-grounded questions}
\PerspectiveTag{\UserIcon}{User}\hfill Context and task performance

\vspace{.4em}
\begin{description}
  \item[Physician] Does the explanation support diagnosis or error detection?
  \item[Engineer] Does it help debug model behavior?
  \item[Auditor] Does it reveal unacceptable dependencies?
  \item[Scientist] Does it generate useful hypotheses?
\end{description}
\end{frame}
